% !TEX root = ../main.tex
\section*{Abstract}
\addcontentsline{toc}{section}{Abstract}

This Master's Thesis presents \texttt{tpeanuts}, a fast, differentiable, GPU-batched framework for neutrino-oscillation phenomenology built on PyTorch and released as public, open-source software, with installation and usage documentation, for the neutrino-phenomenology community. Starting from a validated PyTorch translation of the PEANUTS toolkit's numerical core, three-flavour oscillations, matter effects and solar-neutrino propagation, the project extends substantially: a modular architecture adds atmospheric-neutrino propagation through the Earth, several interchangeable solar and Earth matter/density models, beyond-the-Standard-Model scenarios, and a forward-folding detector layer coupled to a gradient-based statistical-inference engine that connects oscillation probabilities directly to published experimental data. Because its numerical core operates on batched PyTorch tensors, \texttt{tpeanuts} evaluates many oscillation configurations, energies or trajectories in parallel rather than one at a time, scales from CPU to GPU with no code changes, and, being differentiable end to end, connects directly to gradient-based inference and machine-learning applications, making it a common, reusable engine for solar, atmospheric and BSM neutrino studies alike.

The work reviews the theoretical basis of neutrino mixing, vacuum oscillations, the Mikheyev--Smirnov--Wolfenstein matter effect, non-standard interactions, sterile neutrinos and the statistical-inference formalism used throughout the analyses. It then describes the software strategy followed in \texttt{tpeanuts}: identification of the relevant PEANUTS modules, translation of numerical operations to tensor operations, preservation of units and conventions, construction of validation tests against the original implementation and against \texttt{nuSQuIDS}, and preparation of benchmarks for CPU and GPU execution.

The code is organised around a modular architecture in which solar, atmospheric and matter models, new-physics scenarios, detector responses and inference procedures can be selected and combined rather than hard-coded into a single implementation. Dedicated modules support beyond-the-Standard-Model physics, including non-standard interactions and sterile neutrinos. Built on top of this engine, a detector layer folds oscillated fluxes into predicted event counts for several real experiments, and an inference layer fits oscillation and detector parameters to their real, published data via automatic-differentiation-based likelihood minimisation. The most complete case, Daya Bay, is validated against the collaboration's own published results: a near/far ratio fit recovers both $\sin^2 2\theta_{13}$ and $\Delta m^2_{31}$ jointly, within 1--2\% of Daya Bay's published values. This establishes \texttt{tpeanuts} as a working foundation for real, gradient-based neutrino-oscillation data analysis, not only for fast forward simulation.

\subsection*{Keywords}
neutrino oscillations; PEANUTS; PyTorch; differentiable programming; GPU computing; open-source software; MSW effect; solar neutrinos; atmospheric neutrinos; reactor antineutrinos; statistical inference.
